\subsection{Corpus and recognition boundary}
The comparison is pinned to ProveIt commit
\texttt{883da8d55452ea097910ed2509a1eefc833ce292}.  It covers the current
PA research tree and \texttt{docs/incoming}: 100 filed research packets
and six current incoming archives.  Every one of the 13,524 tracked
\texttt{.tm}, \texttt{.json}, and \texttt{.zip} paths selected under those
two roots was read from its pinned Git blob and checked against its Git
content identity.  There are zero missing selected blobs.  Nested archives
were inspected recursively, with identical archive payloads deduplicated.

The recognizer accepts complete numeric five-column tables and four
explicit JSON representations: whole multiline table strings, lists of
five-field rows, lists of row strings, and state/read-keyed dictionaries
of three-field actions.  It does not execute generators or infer machines
from prose, Lean source, compressed opaque objects, or unrecognized
encodings.  Within this stated recognition scope the inventory has 7,328
table occurrences, 1,380 distinct normalized labeled tables, and 1,164
spanning transition-graph classes across all state counts.

The current 194-control comparison contains 548 spanning graph classes
and two nonspanning diagnostic literal forms.  Every one of the 528 older
194-state classes in the maintained packet catalogue has a byte-identical
representative in the fresh census.  All 55 older 193-state catalogue
classes are likewise covered; with the latest incoming experiments the
fresh census has 61 spanning 193-control classes.  These numbers include
rejected candidates, runtime cores and experimental tables.  None is
promoted to a consistency machine solely from its number of rows or a
bounded execution.  The source-attributed roster separately identifies
104 distinct existing named positive-claim graphs.  Each is bound to a
freshly parsed literal graph and a recorded mechanism description.

Exactly 22 additional \texttt{.tm} fragments fail the unique external-halt
condition.  A separate audit finds that they have between two and 139
controls, with either no external halt or multiple unresolved external
targets; none has 194 controls.  Seven recognized tables have no spanning
root, including the two 194-control diagnostics.  Their lack of a spanning
root is invariant under control renaming, reflection and symbol exchange,
so they cannot be isomorphic to any of the four new spanning constructions.

\subsection{Comparison equivalence and restricted negative result}
For a spanning table, an ordered breadth-first traversal reads the zero
transition before the one transition and renumbers controls by discovery
order.  The lexicographically least resulting action tuple over all
spanning roots is a canonical form for the unpointed labeled transition
graph.  The novelty test further minimizes over tape reflection and global
symbol exchange and erases write/move decorations on transitions into
halt.  Thus absence of a match is stronger than nonidentity under an
initial-state-preserving renaming.  Global bit exchange is used only for
structural comparison, not as a claimed semantic equivalence on blank zero
tape.  Exact lexicographic prefix pruning speeds the computation without
changing the minimum.  The offline verifier recomputes all canonical forms
and checks transformed positive controls.

All four new tables are absent from this entire recognized corpus and
are mutually different under that broader comparison.  The same verifier
checks the greatest immediate-action congruence of every one of the 548
existing 194-control spanning graphs.  Each has 194 congruence classes.
This excludes an ordinary action-preserving quotient of those tables;
it is not a lower bound against contextual folds, positive stuttering,
changes of source program, or other machine architectures.

The ledgers retain exact source locations, byte hashes, graph classes,
reported mechanism and proof boundary.  A source report's positive claim
is not an independent audit of its entire PA calibration, and structural
novelty is not by itself a proof of consistency equivalence.  The new
source argument, literal physical certificates and explicit inherited
written calibration dependencies supply those separate obligations.
